% This contents of this file will be inserted into the _Solutions version of the
% output tex document.  Here's an example:

% If assignment with subquestion (1.a) requires a written response, you will
% find the following flag within this document: <SCPD_SUBMISSION_TAG>_1a
% In this example, you would insert the LaTeX for your solution to (1.a) between
% the <SCPD_SUBMISSION_TAG>_1a flags.  If you also constrain your answer between the
% START_CODE_HERE and END_CODE_HERE flags, your LaTeX will be styled as a
% solution within the final document.

% Please do not use the '<SCPD_SUBMISSION_TAG>' character anywhere within your code.  As expected,
% that will confuse the regular expressions we use to identify your solution.

\def\assignmentnum{2 }
\def\assignmenttitle{XCS234 Assignment \assignmentnum}
\input{macros}
\begin{document}
\pagestyle{myheadings} \markboth{}{\assignmenttitle}

% <SCPD_SUBMISSION_TAG>_entire_submission

This handout includes space for every question that requires a written response.
Please feel free to use it to handwrite your solutions (legibly, please).  If
you choose to typeset your solutions, the |README.md| for this assignment includes
instructions to regenerate this handout with your typeset \LaTeX{} solutions.
\ruleskip

\LARGE
2.a
\normalsize

% <SCPD_SUBMISSION_TAG>_2a
\begin{answer}

  \begin{table}[h]
  \centering
  \begin{tabular}{cc}
      \toprule
      Parameter & Value \\
      \midrule
      $w_0$ & \textbf{
        % ### START CODE HERE ###
        1
        % ### END CODE HERE ###
      } \\
      $w_1$ & \textbf{
        % ### START CODE HERE ###
        1
        % ### END CODE HERE ###
      } \\
      $w_2$ & \textbf{
        % ### START CODE HERE ###
        1
        % ### END CODE HERE ###
      } \\
      $w_3$ & \textbf{
        % ### START CODE HERE ###
        -1
        % ### END CODE HERE ###
      } \\
      \bottomrule
  \end{tabular}
  \label{table:Problem_1}
  \caption{Parameters of your reward function.}
  \end{table}
  
\end{answer}
% <SCPD_SUBMISSION_TAG>_2a
\clearpage

\LARGE
2.bi
\normalsize

% <SCPD_SUBMISSION_TAG>_2bi
\begin{answer}
  % ### START CODE HERE ###
  The agent firmly grasps the cube against the palm, holds it nearly still, and slowly rotates it toward the goal orientation.
  This behavior arises because the large position and velocity penalties ($w_0=20$, $w_2=10$) incentivize keeping the cube close and stationary, while the small actuator penalty ($w_3=0.1$) allows the agent to apply the large forces needed for a stable grasp.
  % ### END CODE HERE ###
\end{answer}
% <SCPD_SUBMISSION_TAG>_2bi
\clearpage

\LARGE
2.bii
\normalsize

% <SCPD_SUBMISSION_TAG>_2bii
\begin{answer}
  % ### START CODE HERE ###
  The agent makes weak, tentative hand movements that fail to maintain a firm grasp on the cube.
  The increased actuator penalty ($w_3=1$) discourages applying the large forces needed for effective manipulation, and the very short planning horizon (0.25) prevents the planner from coordinating multi-step actions.
  % ### END CODE HERE ###
\end{answer}
% <SCPD_SUBMISSION_TAG>_2bii
\clearpage

\LARGE
2.biii
\normalsize

% <SCPD_SUBMISSION_TAG>_2biii
\begin{answer}
  % ### START CODE HERE ###
  The agent applies no force and remains still, causing the cube to fall out of the hand and immediately terminating the episode.
  Since only actuator effort is penalized and there is no incentive to hold or orient the cube ($w_0=w_1=w_2=0$), the reward-maximizing behavior is to do nothing.
  % ### END CODE HERE ###
\end{answer}
% <SCPD_SUBMISSION_TAG>_2biii
\clearpage


\LARGE
3.cii
\normalsize

% <SCPD_SUBMISSION_TAG>_3cii
\begin{answer}
  % ### START CODE HERE ###
  % ### END CODE HERE ###
\end{answer}
% <SCPD_SUBMISSION_TAG>_3cii
\clearpage

\LARGE
3.civ
\normalsize

% <SCPD_SUBMISSION_TAG>_3civ
\begin{answer}
  % ### START CODE HERE ###
  % ### END CODE HERE ###
\end{answer}
% <SCPD_SUBMISSION_TAG>_3civ
\clearpage

% <SCPD_SUBMISSION_TAG>_entire_submission

\end{document}
